% !TEX root = ../../../main.tex
% Block card: detection.sax_dsax -- "Symbolic encoding (dSAX)" (Part I, the blocks chapter).
\cardfield{What it does}
Cuts a span into equal pieces and gives each piece a letter for how much the signal rose or
fell across it: with the default three letters, \emph{a} for down, \emph{b} for about the
same, \emph{c} for up. (The SAX family is explained on the card for
\blockref{detection.sax_csax}; dSAX codes a piece's \emph{rise} where the others code its
average.)

\cardfield{When to use it}
When you are after shapes rather than levels: a sharp rise then a long fall becomes a
pattern of letters such as \emph{cc} followed by several \emph{a}s, which
\blockref{detection.symbol_search} can find anywhere in the string. This is the route
research question RQ4 uses: band-pass, remove the baseline, dSAX, then symbol search.

\cardfield{What goes in, what comes out}
In goes a span of one channel, ideally with its baseline removed. Out comes a string with
one letter per piece. A piece that touches a missing sample gets a dash, never a letter.

\cardfield{Settings that matter}
The piece length works as in cSAX (default 20\,s per letter, with the same live read-outs on
the block page). \emph{threshold\_mode} decides where ``about the same'' ends:
\emph{learned} (the default) fits the boundaries to the rises seen in this span;
\emph{absolute} uses a rise you state, in the recording's own unit per piece;
\emph{quantile} makes a share you choose, \emph{same\_fraction}, come out as ``same''.
\emph{trend\_estimator} decides how a rise is measured; the default fits a straight line
through each piece, which is steadier under noise than comparing its two ends.

\cardfield{Watch out for}
In \emph{learned} mode the letters are balanced even on pure noise: about half the pieces
come out as rises or falls when nothing is happening. Do not read a learned string as
structure until you have compared it with a null. A slow drift pushes every piece towards
\emph{c} or \emph{a}; remove the baseline first. Keep the alphabet odd: an even one has no
``same'' letter at all. The app spells the letters \emph{a, b, c}, not D, S, U.
